% Ad Familiares 10.6 (ad-familiares-10-06) --- English translation
% Translated by Claude (Anthropic) from the Latin (Perseus).
% Style guide: STYLE.md.

\ciceroLetterOpener{CICERO PLANCO.}

\ciceroSection{1} What our friend Furnius reported about your feeling toward the state was most gratifying to the Senate and most welcome to the Roman people; but the letter from you that was read out in the Senate seemed in no way to agree with what Furnius had said. You presented yourself as an advocate of peace --- at a moment when your colleague, that most illustrious man, is being besieged by the foulest band of brigands, men who ought either to lay down their arms and ask for peace or, if they keep fighting while demanding it, must have peace wrung from them by victory and not by treaty. As to how the letter on peace, whether Lepidus's or yours, has been received, you can learn from your excellent brother and from G. Furnius.

\ciceroSection{2} My affection for you has moved me, even though your own judgement leaves nothing to be desired and the goodwill of your brother and Furnius and his loyal good sense will be at hand to support you, to wish nonetheless that some lesson of my authority too, for the sake of our many bonds of friendship, should reach you. Believe me, then, Plancus: all the steps of distinction which you have so far attained --- and you have attained the most illustrious --- will carry the names of honours, not the marks of dignity, unless you bind yourself to the liberty of the Roman people and to the authority of the Senate. Detach yourself, I beg you, at long last, from those with whom no judgement of your own but the bonds of the times have linked you.

\ciceroSection{3} In times of public confusion several men have been called consul of whom none was reckoned a consular except the one who showed himself, in spirit, a consular for the state. Such, then, you ought to be: first detach yourself from the company of disloyal citizens who in no way resemble you; next, offer yourself to the Senate and to all loyal men as their adviser, their leader, their captain; and last, judge that peace lies not in arms laid down but in the laying down of the fear of arms and of slavery. If you both act on these principles and feel them, then you will be not merely a consul and a consular, but a great consul and a great consular too; but if otherwise, then in those most illustrious titles of office there will be not merely no dignity but the utmost deformity. I have written somewhat sternly, driven by my affection; you will recognize the truth of it, by putting it to the test in the way that is worthy of you. Dispatched the 20th of March.
